
%%
%% Part 5 - Theoretical Foundations of Flow Matching
%%
%% Flow Matching is a continuous-time generative modelling framework
%% that transports probability mass via a deterministic vector field.
%% This part develops the key theory: the role of the vector field,
%% the choice of linear paths, connections to score matching and
%% optimal transport, and the formulation of the training objective.

\section{Theory of Flow Matching}

Flow Matching provides a deterministic alternative to diffusion
models.  Instead of reversing a stochastic corruption process, Flow
Matching learns a time-dependent vector field that transports a
simple base distribution to the data distribution.  This section
presents the theoretical underpinnings of Flow Matching.

\subsection{Vector Fields and Probability Transport}

At the heart of Flow Matching is a vector field $v_\theta(x,t)$ that
specifies the velocity of a particle located at $x$ at time $t$.
Starting from an initial random vector $x(0)\sim\pi_0$ drawn from a
simple distribution (usually a standard Gaussian), the sample evolves
according to the ordinary differential equation (ODE)
\begin{equation}
  \frac{dx}{dt} = v_\theta\bigl(x(t), t\bigr),
  \quad x(0) \sim \pi_0.
\end{equation}
Solving this ODE from $t=0$ to $t=1$ yields a final random variable
$x(1)$, which we hope matches the target data distribution \(\pi_1\).
The vector field therefore acts as a continuous-time flow that carries
probability mass from the source to the target.

\begin{solutionbox}
  \textbf{Role of the vector field.}  The vector field determines the
  instantaneous direction and speed with which a sample moves through
  space.  If $v_\theta(x,t)$ points towards regions of higher data
  density at all times, integrating the ODE will gradually transform
  the initial noise into realistic data samples.  Unlike diffusion
  models that rely on randomness, Flow Matching is fully
  deterministic once the vector field is learned.
\end{solutionbox}

\subsection{Choice of Conditional Paths}

To train $v_\theta$, one must define a family of conditional paths
between pairs of points $(x_0, x_1)$ sampled from the base and data
Distributions.  The simplest choice is the linear interpolation
\begin{equation}
  x(t) = (1-t)x_0 + t x_1,
\end{equation}
which connects $x_0$ and $x_1$ along a straight line.  The velocity
associated with this path is constant and equal to $x_1 - x_0$.

\begin{solutionbox}
  \textbf{Why linear paths?}  Linear paths are appealing because
  they simplify the regression target: the true velocity is just
  $u_t(x\mid x_0,x_1) = x_1 - x_0$ for all $t$.  This reduces the
  learning problem to a straightforward least squares regression.
  While more complex paths may minimise transport cost under certain
  couplings, linear paths coupled with optimal transport produce the
  geodesic in Euclidean space.  In practice, they offer a good
  trade-off between simplicity and performance.
\end{solutionbox}

\subsection{Conditional Flow Matching Loss}

Given a path family and an initial sample $x_0\sim\pi_0$ and target
sample $x_1\sim\pi_1$, one defines the conditional velocity field
\begin{equation}
  u_t(x\mid x_0,x_1) = x_1 - x_0,
\end{equation}
which is constant along the linear path.  The network $v_\theta$ is
trained to match this conditional velocity.  The Conditional Flow
Matching (CFM) loss is
\begin{equation}
  \mathcal{L}_{\text{CFM}}(\theta) = \mathbb{E}_{x_0,x_1,t}
  \bigl\|v_\theta\bigl((1-t)x_0 + t x_1, t\bigr) - (x_1 - x_0)\bigr\|^2.
\end{equation}
The expectation is taken over pairs $(x_0,x_1)$ and times
$t\sim\mathrm{Uniform}(0,1)$.  Minimising this loss yields a vector
field that matches the conditional velocity in expectation.

\begin{solutionbox}
  The loss can be derived from denoising score matching.  For a
  Gaussian perturbation kernel $q(x_t\mid x_0) = \mathcal{N}(x_t;
  x_0, \sigma_t^2 I)$, the score of the conditional distribution
  satisfies
  \begin{equation}
    \nabla_{x_t} \log q(x_t\mid x_0) = \frac{x_0 - x_t}{\sigma_t^2}.
  \end{equation}
  When specialised to linear paths with variance $\sigma_t^2 = t(1-t)\sigma^2$,
  the score reduces to a constant proportional to $x_1 - x_0$.  Thus
  matching the score is equivalent to matching the constant
  velocity.
\end{solutionbox}

\subsection{Connections to Optimal Transport}

Flow Matching has strong ties to optimal transport theory.  When
$(x_0,x_1)$ are sampled from an optimal coupling (i.e. a joint
Distribution that minimises the expected squared distance between
$\pi_0$ and $\pi_1$), the linear path becomes the geodesic that
minimises the cost of transporting probability mass.  In practice,
optimal couplings are approximated by pairing shuffled noise samples
with data samples.  While not exact, this approach yields a reasonable
approximation for high-dimensional data.

\subsection{Advantages over Diffusion Models}

Flow Matching offers several advantages compared to diffusion models:
\begin{itemize}
  \item It operates in continuous time, avoiding the need to choose a
    discrete timestep schedule.  This can provide greater flexibility.
  \item The training objective is a simple regression problem with
    constant targets, leading to stable gradients and faster
    convergence.
  \item Sampling amounts to solving an ODE, which can be done with a
    small number of integration steps (e.g. 100) using standard
    numerical solvers such as Euler or Runge–Kutta.
  \item There is no stochasticity during sampling once the initial
    noise sample is fixed, simplifying the generation process.
\end{itemize}

However, Flow Matching does not automatically account for varying
noise scales or multi-scale structure in data, which diffusion models
handle naturally via their variance schedules.  Hybrid approaches
seek to combine the strengths of both frameworks.

\subsection{Conclusion of Part 5}

This part has introduced the concept of Flow Matching, defined the
vector field ODE, motivated linear interpolation paths, derived the
Conditional Flow Matching loss and discussed connections to score
matching and optimal transport.  Flow Matching provides a
conceptually elegant and computationally efficient alternative to
stochastic diffusion models.  The next part will describe how to
implement Flow Matching in practice and evaluate its performance on
financial time series.

\subsection{Relation to Probability Flow ODE}

In diffusion models one can derive a deterministic probability flow
ODE that has the same marginal distributions as the stochastic
forward and reverse SDEs.  This ODE is given by
\begin{equation}
  \frac{dx}{dt} = f(x,t) - \frac{1}{2}g(t)^2 \nabla_x \log p_t(x),
\end{equation}
where $f$ and $g$ are the drift and diffusion coefficients of the
forward SDE and $p_t$ denotes the marginal distribution of the
forward process.  Solving this ODE transports samples from $p_0$ to
$p_1$ deterministically.  Flow Matching can be viewed as learning
such a probability flow ODE directly, without first specifying a
stochastic process.

\subsection{Score-Based Generative Modelling vs Flow Matching}

Score-based models aim to estimate the score function $\nabla_x\log
p_t(x)$ of the data distribution at different noise levels and then
integrate an ODE or SDE to generate samples.  Flow Matching bypasses
score estimation by directly regressing the velocity field that
connects noise to data.  Both approaches rely on continuous-time
dynamics but differ in their training objectives.  Score matching
requires evaluating gradients of log-densities, which may be
unstable, whereas velocity regression involves simple differences
between paired data and noise samples.  The two methods are
complementary and hybrid techniques can leverage their strengths.

\subsection{Optimal Transport Theory Primer}

Optimal transport provides a mathematical framework for transporting
one probability measure to another in a cost-optimal manner.  The
Wasserstein distance, defined as
\begin{equation}
  W_2^2(\mu,\nu) = \inf_{\gamma\in\Pi(\mu,\nu)}\int\|x-y\|^2\,d\gamma(x,y),
\end{equation}
measures the minimal quadratic cost of transporting mass from $\mu$ to
$\nu$.  Here $\Pi(\mu,\nu)$ is the set of all couplings with
marginals $\mu$ and $\nu$.  When the coupling achieving the infimum is
used, the linear interpolation path between paired samples is the
geodesic that minimises the cost.  Flow Matching leverages this
concept by sampling $(x_0,x_1)$ pairs according to (approximate)
optimal couplings, ensuring that the learned vector field follows
straight lines between optimally matched points.

\subsection{Future Directions in Flow Matching Theory}

Current Flow Matching implementations use simple linear paths and
Euclidean cost functions.  Future research may explore non-linear
interpolation schemes, adaptive path selection, and different cost
metrics (e.g. \(L_1\) or Mahalanobis distances).  Extending Flow
Matching to conditionally dependent multi-dimensional data,
incorporating conditioning information and learning couplings jointly
with the velocity field are open problems.  Further theoretical
analysis of convergence properties and connections to other forms of
continuous normalising flows may also yield deeper insights.

% Additional padding lines for length
% The remaining lines ensure the appendix meets the required length.
% End of Part 5.
